\documentclass[12pt,a4paper]{article}

\usepackage[left=30mm,top=20mm,right=15mm,bottom=33mm,includeheadfoot]{geometry}
\usepackage[utf8]{inputenc}
\usepackage[T2A]{fontenc}
\usepackage[russian,english]{babel}
\usepackage{setspace}
\usepackage{indentfirst}
\usepackage{array}
\usepackage{tabularx}
\usepackage{longtable}
\usepackage{multirow}
\usepackage{makecell}
\usepackage{graphicx}
\usepackage{hyperref}
\usepackage{xurl}
\usepackage{fancyhdr}
\usepackage{titlesec}
\usepackage{enumitem}
\usepackage{lastpage}
\usepackage[absolute,overlay]{textpos}

\hypersetup{hidelinks}
\onehalfspacing
\setlength{\parindent}{1.25cm}
\setlength{\headheight}{30pt}
\setlength{\footskip}{72pt}
\emergencystretch=3em
\setlength{\TPHorizModule}{1mm}
\setlength{\TPVertModule}{1mm}
\textblockorigin{0mm}{0mm}

\newcommand{\doccodebase}{RU.17701729.05.05-01 12 01-1}
\newcommand{\doccode}{RU.17701729.05.05-01 12 01-1}
\newcommand{\doccodeapproval}{RU.17701729.05.05-01 12 01-1-ЛУ}
\newcommand{\doctitle}{Андроид-приложение для совместного планирования поездок}
\newcommand{\doctype}{Текст программы}

\titleformat{\section}{\large\bfseries\centering}{\thesection}{1em}{\MakeUppercase}
\titleformat{\subsection}{\normalsize\bfseries}{\thesubsection}{1em}{}
\titleformat{\subsubsection}{\normalsize\bfseries}{\thesubsubsection}{1em}{}
\titlespacing*{\section}{0pt}{1.5ex plus 0.5ex minus 0.3ex}{1.0ex}
\titlespacing*{\subsection}{0pt}{1.2ex plus 0.4ex minus 0.2ex}{0.8ex}
\titlespacing*{\subsubsection}{0pt}{1.0ex plus 0.3ex minus 0.2ex}{0.5ex}

\setlist[enumerate]{leftmargin=1.2cm,itemsep=0.2em,topsep=0.3em}
\setlist[itemize]{leftmargin=1.2cm,itemsep=0.2em,topsep=0.3em}

\newcommand{\pagefooter}{%
  \begingroup
  \setlength{\tabcolsep}{3pt}%
  \renewcommand{\arraystretch}{1.0}%
  \footnotesize
  \begin{tabularx}{\textwidth}{|>{\centering\arraybackslash}p{0.305\textwidth}|>{\centering\arraybackslash}p{0.14\textwidth}|>{\centering\arraybackslash}p{0.13\textwidth}|>{\centering\arraybackslash}p{0.125\textwidth}|>{\centering\arraybackslash}X|}
    \hline
    Изм. & Лист & № докум. & Подп. & Дата \\
    \hline
    \doccode & & & & \\
    \hline
    Инв.~№~подл. & Подп.~и~дата & Взам.~инв.~№ & Инв.~№~дубл. & Подп.~и~дата \\
    \hline
  \end{tabularx}
  \endgroup
}

\fancypagestyle{tpstyle}{
  \fancyhf{}
  \fancyhead[C]{\textbf{\thepage}\\[-0.2ex]\textbf{\doccode}}
  \fancyfoot[C]{\pagefooter}
  \renewcommand{\headrulewidth}{0pt}
  \renewcommand{\footrulewidth}{0pt}
}

\fancypagestyle{plain}{
  \fancyhf{}
  \fancyhead[C]{\textbf{\thepage}\\[-0.2ex]\textbf{\doccode}}
  \fancyfoot[C]{\pagefooter}
  \renewcommand{\headrulewidth}{0pt}
  \renewcommand{\footrulewidth}{0pt}
}

\newcommand{\leftapprovalstamp}[1]{%
  \begingroup
  \renewcommand{\arraystretch}{1.0}
  \setlength{\tabcolsep}{0pt}
  \footnotesize
  \begin{tabular}{|>{\centering\arraybackslash}m{0.65cm}|>{\centering\arraybackslash}m{1.05cm}|}
    \hline
    \rotatebox{90}{Подп. и дата} & \rule{0pt}{1.10cm}\\
    \hline
    \rotatebox{90}{Инв. № дубл.} & \rule{0pt}{1.00cm}\\
    \hline
    \rotatebox{90}{Взам. инв. №} & \rule{0pt}{1.00cm}\\
    \hline
    \rotatebox{90}{Подп. и дата} & \rule{0pt}{1.10cm}\\
    \hline
    \rotatebox{90}{Инв. № подл.} &
    \rotatebox{90}{\scalebox{0.70}{#1}}\rule{0pt}{0.95cm}\\
    \hline
  \end{tabular}
  \endgroup
}

\newcommand{\placestamp}[2]{%
  \begin{textblock*}{19mm}(1mm,#2)
    \leftapprovalstamp{#1}
  \end{textblock*}
}

\newcommand{\registrationblankrow}{%
  \rule{0pt}{0.78cm} & & & & & & & & & \\
  \hline
}

\newcommand{\registrationtable}{%
  \begingroup
  \setlength{\tabcolsep}{2pt}%
  \renewcommand{\arraystretch}{1.0}
  \footnotesize
  \centering
  \begin{tabular}{|>{\centering\arraybackslash}p{0.75cm}|>{\centering\arraybackslash}p{1.55cm}|>{\centering\arraybackslash}p{1.55cm}|>{\centering\arraybackslash}p{1.55cm}|>{\centering\arraybackslash}p{1.7cm}|>{\centering\arraybackslash}p{1.9cm}|>{\centering\arraybackslash}p{2.2cm}|>{\centering\arraybackslash}p{2.4cm}|>{\centering\arraybackslash}p{1.0cm}|>{\centering\arraybackslash}p{1.0cm}|}
    \hline
    \multicolumn{10}{|c|}{\textbf{Лист регистрации изменений}}\\
    \hline
    \multicolumn{5}{|c|}{\makecell[c]{Номера листов\\(страниц)}} & & & & & \\
    \hline
    Изм. & \makecell{Изме-\\ненных} & \makecell{Заме-\\ненных} & Новых & \makecell{Аннулиро-\\ванных} &
    \makecell[c]{Всего\\листов\\(страниц)\\в докум.} &
    \makecell[c]{№\\документа} &
    \makecell[c]{Входящий\\№\\сопроводит.\\документа\\и\\дата} &
    \makecell[c]{Подп.} &
    \makecell[c]{Дата} \\
    \hline
    \registrationblankrow
    \registrationblankrow
    \registrationblankrow
    \registrationblankrow
    \registrationblankrow
    \registrationblankrow
    \registrationblankrow
    \registrationblankrow
    \registrationblankrow
    \registrationblankrow
    \registrationblankrow
    \registrationblankrow
    \registrationblankrow
    \registrationblankrow
  \end{tabular}
  \endgroup
}

\begin{document}
\hypersetup{pageanchor=false}
\selectlanguage{russian}

\thispagestyle{empty}
\placestamp{\doccode}{146mm}
\vspace*{-2.1cm}

\begin{center}
  {\small\textbf{ПРАВИТЕЛЬСТВО РОССИЙСКОЙ ФЕДЕРАЦИИ}}\\
  {\small\textbf{ФЕДЕРАЛЬНОЕ ГОСУДАРСТВЕННОЕ АВТОНОМНОЕ}}\\
  {\small\textbf{ОБРАЗОВАТЕЛЬНОЕ УЧРЕЖДЕНИЕ ВЫСШЕГО ОБРАЗОВАНИЯ}}\\
  {\small\textbf{НАЦИОНАЛЬНЫЙ ИССЛЕДОВАТЕЛЬСКИЙ УНИВЕРСИТЕТ}}\\
  {\small\textbf{«ВЫСШАЯ ШКОЛА ЭКОНОМИКИ»}}\\[0.35cm]
  Факультет компьютерных наук\\
  Образовательная программа «Программная инженерия»
\end{center}

\vspace{0.6cm}
\noindent
\begin{minipage}[t]{0.49\textwidth}
  \centering
  \textbf{СОГЛАСОВАНО}\\
  Руководитель,\\
  Доцент департамента «Программной инженерии»\\[0.55cm]
  \underline{\hspace{2.7cm}} А.Д. Брейман\\
  «\underline{\hspace{0.9cm}}» \underline{\hspace{2.1cm}} 2026 г.
\end{minipage}\hfill
\begin{minipage}[t]{0.49\textwidth}
  \centering
  \textbf{УТВЕРЖДАЮ}\\
  Академический руководитель\\
  образовательной программы\\
  «Программная инженерия»,\\
  старший преподаватель департамента\\
  программной инженерии\\[0.2cm]
  \underline{\hspace{2.7cm}} Н.А. Павлочев\\
  «\underline{\hspace{0.9cm}}» \underline{\hspace{2.1cm}} 2026 г.
\end{minipage}

\vspace{0.75cm}
\begin{center}
  {\large\textbf{\doctitle}}\\[0.22cm]
  {\large\textbf{\doctype}}\\[0.18cm]
  {\large\textbf{ЛИСТ УТВЕРЖДЕНИЯ}}\\[0.18cm]
  {\large\textbf{\doccodeapproval}}
\end{center}

\vspace{1.15cm}
\begin{flushright}
  Исполнитель:\\[0.25cm]
  студент группы БПИ223\\[0.35cm]
  \underline{\hspace{2.5cm}} И.И. Новик\\[0.25cm]
  «\underline{\hspace{0.9cm}}» \underline{\hspace{2.8cm}} 2026 г.
\end{flushright}

\begin{textblock*}{\paperwidth}(0mm,281mm)
  \centering{\large\textbf{2026}}
\end{textblock*}
\newpage

\thispagestyle{empty}
\placestamp{\doccodeapproval}{144mm}
\vspace*{-0.6cm}

\begin{flushleft}
  \textbf{УТВЕРЖДЕН}\\[0.2cm]
  \textbf{\doccodeapproval}
\end{flushleft}

\vspace{1.4cm}
\begin{center}
  {\large\textbf{\doctitle}}\\[0.22cm]
  {\Large\textbf{\doctype}}\\[0.22cm]
  {\Large\textbf{\doccode}}\\[0.2cm]
  {\Large\textbf{Листов \pageref{lastnumberedpage}}}
\end{center}

\begin{textblock*}{\paperwidth}(0mm,281mm)
  \centering{\large\textbf{2026}}
\end{textblock*}

\clearpage
\hypersetup{pageanchor=true}
\setcounter{page}{2}
\pagestyle{tpstyle}
\tableofcontents
\clearpage

\section*{ТЕКСТ ПРОГРАММЫ}
\addcontentsline{toc}{section}{ТЕКСТ ПРОГРАММЫ}

\subsection*{1 Клиентская часть}
Клиентская часть CoTrip реализована в каталоге \path{android/app/src/main/java/nvk/cotrip}.
Исходный код опубликован в репозитории \url{https://github.com/ilya-nvk/cotrip}.

\subsubsection*{1.1 Структура клиентской части}
\begin{enumerate}
  \item \texttt{ui} --- экраны и пользовательские сценарии (авторизация, поездки, идеи, маршрут, расходы, настройки, погода, рекомендации ИИ).
  \item \texttt{data/repository} --- прикладная логика сценариев и координация источников данных.
  \item \texttt{data/network} --- API-клиенты, перехватчики, сетевые модели.
  \item \texttt{data/cache} --- локальные кэш-хранилища клиентских сущностей.
  \item \texttt{data/sync} --- офлайн-очередь и синхронизация после восстановления сети.
  \item \texttt{notifications} --- обработка уведомлений и переходов по событиям.
  \item \texttt{di} --- конфигурация внедрения зависимостей.
\end{enumerate}

\subsubsection*{1.2 Основные библиотеки клиентской части}
\begin{enumerate}
  \item Kotlin.
  \item Jetpack Compose.
  \item AndroidX Navigation.
  \item Hilt/Dagger.
  \item Retrofit + OkHttp + Kotlinx Serialization.
  \item Room.
  \item WorkManager.
  \item Firebase Cloud Messaging SDK.
\end{enumerate}

\subsubsection*{1.3 Ссылки на исходный код клиентской части}
\begin{enumerate}
  \item Общий репозиторий: \url{https://github.com/ilya-nvk/cotrip}
  \item Клиентский код: \url{https://github.com/ilya-nvk/cotrip/tree/main/android/app/src/main/java/nvk/cotrip}
\end{enumerate}

\subsection*{2 Серверная часть}
Серверная часть CoTrip реализована в каталоге \path{backend/src/main/kotlin/nvk/cotrip/backend}.

\subsubsection*{2.1 Структура серверной части}
\begin{enumerate}
  \item \texttt{routes} --- HTTP API-маршруты (\texttt{auth}, \texttt{trips}, \texttt{ideas}, \texttt{itinerary}, \texttt{expenses}, \texttt{sync}, \texttt{notifications}).
  \item \texttt{db} --- доступ к данным и репозитории предметных сущностей.
  \item \texttt{auth} --- валидация токенов и управление сессиями.
  \item \texttt{ws} --- WebSocket-канал для обсуждений идей.
  \item \texttt{integrations} --- внешние интеграции (погода, AI).
  \item \texttt{notifications} --- серверная публикация push-событий.
  \item \texttt{plugins} --- конфигурация Ktor (сериализация, маршрутизация, обработка ошибок, безопасность).
\end{enumerate}

\subsubsection*{2.2 Основные библиотеки серверной части}
\begin{enumerate}
  \item Kotlin.
  \item Ktor.
  \item PostgreSQL JDBC.
  \item Kotlinx Serialization.
  \item Firebase Admin SDK.
\end{enumerate}

\subsubsection*{2.3 Ссылки на исходный код серверной части}
\begin{enumerate}
  \item Серверный код: \url{https://github.com/ilya-nvk/cotrip/tree/main/backend/src/main/kotlin/nvk/cotrip/backend}
  \item API-маршруты: \url{https://github.com/ilya-nvk/cotrip/tree/main/backend/src/main/kotlin/nvk/cotrip/backend/routes}
\end{enumerate}

\clearpage
\section*{СПИСОК ИСТОЧНИКОВ}
\addcontentsline{toc}{section}{СПИСОК ИСТОЧНИКОВ}
\begin{enumerate}
  \item ГОСТ 19.101-77: Виды программ и программных документов. // Единая система программной документации. --- М.: ИПК Издательство стандартов, 2001.
  \item ГОСТ 19.102-77: Стадии разработки. // Единая система программной документации. --- М.: ИПК Издательство стандартов, 2001.
  \item ГОСТ 19.103-77: Обозначения программ и программных документов. // Единая система программной документации. --- М.: ИПК Издательство стандартов, 2001.
  \item ГОСТ 19.104-78: Основные надписи. // Единая система программной документации. --- М.: ИПК Издательство стандартов, 2001.
  \item ГОСТ 19.105-78: Общие требования к программным документам. // Единая система программной документации. --- М.: ИПК Издательство стандартов, 2001.
  \item ГОСТ 19.106-78: Требования к программным документам, выполненным печатным способом. // Единая система программной документации. --- М.: ИПК Издательство стандартов, 2001.
  \item ГОСТ 19.602-78: Правила дублирования, учета и хранения программных документов, выполненных печатным способом. // Единая система программной документации. --- М.: ИПК Издательство стандартов, 2001.
  \item ГОСТ 19.603-78: Общие правила внесения изменений. // Единая система программной документации. --- М.: ИПК Издательство стандартов, 2001.
  \item GitHub-репозиторий проекта CoTrip. [Электронный ресурс] --- URL: \url{https://github.com/ilya-nvk/cotrip} режим доступа: свободный (дата обращения: 19.03.2026).
  \item GitHub-репозиторий с текстом клиентской части. [Электронный ресурс] --- URL: \url{https://github.com/ilya-nvk/cotrip/tree/main/android/app/src/main/java/nvk/cotrip} режим доступа: свободный (дата обращения: 19.03.2026).
  \item GitHub-репозиторий с текстом серверной части. [Электронный ресурс] --- URL: \url{https://github.com/ilya-nvk/cotrip/tree/main/backend/src/main/kotlin/nvk/cotrip/backend} режим доступа: свободный (дата обращения: 19.03.2026).
  \item GitHub-репозиторий с текстом клиентских экранов. [Электронный ресурс] --- URL: \url{https://github.com/ilya-nvk/cotrip/tree/main/android/app/src/main/java/nvk/cotrip/ui} режим доступа: свободный (дата обращения: 19.03.2026).
  \item GitHub-репозиторий с текстом API-маршрутов сервера. [Электронный ресурс] --- URL: \url{https://github.com/ilya-nvk/cotrip/tree/main/backend/src/main/kotlin/nvk/cotrip/backend/routes} режим доступа: свободный (дата обращения: 19.03.2026).
  \item Kotlin Documentation. [Электронный ресурс] --- URL: \url{https://kotlinlang.org/docs/home.html} режим доступа: свободный (дата обращения: 19.03.2026).
  \item Android Developers: Jetpack Compose. [Электронный ресурс] --- URL: \url{https://developer.android.com/compose} режим доступа: свободный (дата обращения: 19.03.2026).
  \item Android Developers: Navigation Compose. [Электронный ресурс] --- URL: \url{https://developer.android.com/develop/ui/compose/navigation} режим доступа: свободный (дата обращения: 19.03.2026).
  \item Android Developers: Hilt. [Электронный ресурс] --- URL: \url{https://developer.android.com/training/dependency-injection/hilt-android} режим доступа: свободный (дата обращения: 19.03.2026).
  \item Dagger Hilt Documentation. [Электронный ресурс] --- URL: \url{https://dagger.dev/hilt/} режим доступа: свободный (дата обращения: 19.03.2026).
  \item Retrofit Documentation. [Электронный ресурс] --- URL: \url{https://square.github.io/retrofit/} режим доступа: свободный (дата обращения: 19.03.2026).
  \item OkHttp Documentation. [Электронный ресурс] --- URL: \url{https://square.github.io/okhttp/} режим доступа: свободный (дата обращения: 19.03.2026).
  \item Kotlinx Serialization Documentation. [Электронный ресурс] --- URL: \url{https://github.com/Kotlin/kotlinx.serialization} режим доступа: свободный (дата обращения: 19.03.2026).
  \item Android Developers: Room. [Электронный ресурс] --- URL: \url{https://developer.android.com/training/data-storage/room} режим доступа: свободный (дата обращения: 19.03.2026).
  \item Android Developers: WorkManager. [Электронный ресурс] --- URL: \url{https://developer.android.com/topic/libraries/architecture/workmanager} режим доступа: свободный (дата обращения: 19.03.2026).
  \item Firebase Cloud Messaging. [Электронный ресурс] --- URL: \url{https://firebase.google.com/docs/cloud-messaging} режим доступа: свободный (дата обращения: 19.03.2026).
  \item Firebase Admin SDK Documentation. [Электронный ресурс] --- URL: \url{https://firebase.google.com/docs/admin/setup} режим доступа: свободный (дата обращения: 19.03.2026).
  \item Ktor Documentation. [Электронный ресурс] --- URL: \url{https://ktor.io/docs/} режим доступа: свободный (дата обращения: 19.03.2026).
  \item PostgreSQL Documentation. [Электронный ресурс] --- URL: \url{https://www.postgresql.org/docs/} режим доступа: свободный (дата обращения: 19.03.2026).
  \item OpenWeather API Documentation. [Электронный ресурс] --- URL: \url{https://openweathermap.org/api} режим доступа: свободный (дата обращения: 19.03.2026).
  \item Yandex Cloud AI Documentation. [Электронный ресурс] --- URL: \url{https://yandex.cloud/ru/docs/foundation-models/} режим доступа: свободный (дата обращения: 19.03.2026).
\end{enumerate}

\clearpage
\thispagestyle{empty}
\phantomsection
\addcontentsline{toc}{section}{Лист регистрации изменений}
\vspace*{-1.1cm}
\begin{center}
  {\Large\textbf{ЛИСТ РЕГИСТРАЦИИ ИЗМЕНЕНИЙ}}
\end{center}
\vspace{0.35cm}
\registrationtable

\label{lastnumberedpage}
\end{document}
